% ============================================================
% CAPITOLO O - METODO COMPLETO PER L'ARMONIZZAZIONE
%              DEL BASSO SENZA NUMERI
% (Capitolo finale - NUOVO)
% ============================================================
\chapter{Metodo operativo per il basso senza numeri}
\label{chap:metodo_basso_senza_numeri}

Questo è il capitolo conclusivo del percorso. Tutto ciò che è stato studiato --- fondamenti acustici, triadi, rivolti, condotta delle voci, funzioni armoniche, cadenze, dissonanze, progressioni, modulazioni --- converge qui in un \textbf{metodo operativo} per l'armonizzazione autonoma di un basso privo di numerica.

Non esiste una formula magica. Esiste però una \textbf{procedura di analisi e scelta} che, se seguita sistematicamente, riduce l'incertezza e permette di produrre realizzazioni musicalmente coerenti. Questo capitolo presenta quella procedura, la illustra su esempi completi, cataloga gli errori tipici e suggerisce strategie di scelta nei casi più ambigui.

\section{Le cinque fasi del metodo}
\label{sec:cinque_fasi}

\begin{regola}[Il metodo in cinque fasi]
	\label{reg:metodo_cinque_fasi}
	L'armonizzazione del basso senza numeri si articola in cinque fasi sequenziali:
	\begin{enumerate}
		\item \textbf{Analisi preliminare} --- Tonalità, modo, struttura metrica.
		\item \textbf{Segmentazione fraseologica} --- Frasi, cadenze, periodo.
		\item \textbf{Classificazione delle note} --- Gradi, funzioni, movimento, alterazioni.
		\item \textbf{Stesura dell'armatura} --- Scrittura delle voci superiori con rispetto delle regole di condotta.
		\item \textbf{Revisione sistematica} --- Controllo degli errori.
	\end{enumerate}
	Queste fasi sono \textit{sequenziali}: non si passa alla fase successiva prima di aver completato la precedente.
\end{regola}

\section{Fase 1 — Analisi preliminare}
\label{sec:fase1_basso}

\subsection{Tonalità e modo}

Prima di toccare carta e penna, dedicare almeno un minuto alla \textbf{lettura silenziosa} dell'intero basso. Porsi le domande:
\begin{itemize}
	\item Qual è l'armatura di chiave?
	\item La nota iniziale è la tonica o la dominante?
	\item La nota finale (di solito lunga) è la tonica?
	\item Ci sono alterazioni cromatiche? Quali e dove?
\end{itemize}

\begin{regola}[Lettura diagnostica del modo]
	\label{reg:lettura_diagnostica}
	\begin{itemize}
		\item Se il basso contiene regolarmente la \textbf{sensibile alzata} (VII grado cromatico), il brano è in modo minore armonico.
		\item Se il basso alterna il VII grado alzato e abbassato, potremmo essere in modo minore melodico: stare attenti al contesto (salita vs.\ discesa).
		\item Se il basso ha solo le note dell'armatura, ma si muove attorno alla tonica di un modo minore, siamo probabilmente in modo minore naturale (raro negli esercizi scolastici).
		\item Un accordo iniziale con la terza maggiore visibile nell'arpeggio del basso conferma il modo maggiore.
	\end{itemize}
\end{regola}

\subsection{Struttura metrica e ritmo armonico previsto}

Leggere il metro (4/4, 3/4, 3/2, ecc.) e stimare il ritmo armonico previsto. In un esercizio scolastico standard:
\begin{itemize}
	\item In \textbf{4/4}: un accordo per ogni semiminima o minima, a seconda delle indicazioni dell'esaminatore.
	\item In \textbf{3/2}: un accordo per minima (tre accordi per misura) o uno per misura intera, a seconda della musica.
	\item La regola generale è: tanti accordi quanti sono gli \textit{accenti principali} del metro.
\end{itemize}

\subsection{Identificazione dei pattern visivi}

Prima ancora di analizzare armonia, cercare \textbf{pattern visibili} nel profilo del basso:
\begin{itemize}
	\item Movimenti scalari lunghi (più di 3-4 note in senso congiunto) $\rightarrow$ Regola dell'Ottava o progressione.
	\item Salti ripetuti in senso discendente $\rightarrow$ progressione per quinte discendenti.
	\item Alternanza semitoni/toni $\rightarrow$ progressione cromatica.
	\item Salti di quarta o quinta seguiti da grado congiunto $\rightarrow$ inizio di progressione per terze.
\end{itemize}

\section{Fase 2 — Segmentazione fraseologica}
\label{sec:fase2_basso}

\subsection{Identificare le cadenze}

Le cadenze si trovano quasi sempre ai \textbf{punti di articolazione ritmica del brano}: fine di una frase di 4 misure, metà del periodo, fine del brano.

\begin{regola}[Segnali di cadenza nel basso senza numeri]
	\label{reg:segnali_cadenza}
	\begin{itemize}
		\item Una nota lunga (semibreve o minima in 4/4) dopo una frase in movimento $\rightarrow$ cadenza.
		\item Il basso che arriva alla tonica dopo un movimento di quarta ascendente o quinta discendente $\rightarrow$ cadenza perfetta.
		\item Il basso che si ferma sulla dominante $\rightarrow$ cadenza sospesa.
		\item Presenza di due note del basso V-I con ritmo regolare prima di una nota lunga $\rightarrow$ cadenza composta.
		\item Una nota alterata seguita immediatamente da un grado congiunto $\rightarrow$ possibile cadenza con sensibile o con tonicizzazione.
	\end{itemize}
\end{regola}

\subsection{Struttura antecedente-conseguente}

In un basso di 8 misure, applicare il modello del periodo:
\begin{itemize}
	\item \textbf{Misure 1-4} (antecedente): la cadenza di metà deve essere \textit{aperta} --- sospesa (V) o imperfetta (I non fondamentale, o soprano non all'ottava).
	\item \textbf{Misure 5-8} (conseguente): la cadenza finale deve essere \textit{chiusa} --- cadenza perfetta (V-I in stato fondamentale, soprano all'ottava o alla quinta).
\end{itemize}

Se il basso è di 4 misure, applicare lo stesso criterio: cadenza di metà aperta, cadenza finale chiusa.

\section{Fase 3 — Classificazione delle note}
\label{sec:fase3_basso}

Questa è la fase più impegnativa. Si passa nota per nota, annotando (idealmente sopra o sotto il basso) la funzione armonica dedotta.

\subsection{Schema di classificazione}

Per ogni nota del basso, rispondere nell'ordine:

\begin{enumerate}
	\item \textbf{Qual è il suo grado nella scala?} (I, II, III\ldots)
	\item \textbf{È su un tempo forte o debole?}
	\item \textbf{Come si muove rispetto alla nota precedente?} (salto, grado congiunto, nota tenuta)
	\item \textbf{Come si muove verso la nota seguente?} (uguale a sopra)
	\item \textbf{Appartiene a un pattern riconoscibile?} (cadenza, progressione, Regola dell'Ottava)
	\item \textbf{Quale accordo è più coerente con il contesto?}
\end{enumerate}

\subsection{Tabella di riferimento rapido}

\begin{center}
\begin{tabular}{|l|l|l|}
\hline
\textbf{Situazione} & \textbf{Soluzione di default} & \textbf{Note} \\
\hline
I grado, tempo forte & I (stato fond.) & Tonica, 5 \\
II grado, congiunto & VII6 o V4/3 & Regola dell'Ottava \\
III grado, congiunto & I6 & Regola dell'Ottava \\
IV grado, congiunto & II6 o II6/5 & Regola dell'Ottava \\
V grado, congiunto & V o V7 & o I6/4 cadenzale/passaggio \\
VI grado, congiunto & IV6 & Regola dell'Ottava \\
VII grado, congiunto & V6 o V6/5 & Regola dell'Ottava \\
Nota con salto & Stato fond. & Evitare rivolti \\
Nota cromatica alt. & Dom. secondaria & Analizzare contesto \\
Nota finale & I (stato fond.) & Cadenza perfetta \\
\hline
\end{tabular}
\end{center}

\subsection{Gestione delle tonicizzazioni nel basso senza numeri}

Quando nel basso appare un'alterazione cromatica non prevista dall'armatura, il primo passo è verificare se si tratta di una \textbf{dominante secondaria} (tonicizzazione).

\begin{regola}[Riconoscimento della dominante secondaria nel basso]
	\label{reg:dom_secondaria_basso}
	\begin{itemize}
		\item Il \Fadiesis{} nel basso di Do maggiore, in corrispondenza del II grado, indica una \textbf{dominante della dominante} (V/V): l'accordo sul Re dovrà essere Re maggiore con \Fadiesis{}.
		\item Il \Sibemolle{} nel basso di Do maggiore, sul I grado, indica una \textbf{dominante della sottodominante} (V/IV): si usa Do con settima minore (\Sibemolle{}).
		\item \Midiesis{} o \Ladiesis{} nel basso in La minore: verificare se sono la sensibile del relativo maggiore (Do) o di un altro grado tonicizzato.
	\end{itemize}
\end{regola}

\section{Fase 4 — Stesura delle voci superiori}
\label{sec:fase4_basso}

Una volta stabilita la ``numerica implicita'' (il piano armonico), si scrivono le voci superiori con le stesse regole usate per il basso figurato.

\begin{regola}[Priorità nella stesura delle voci superiori]
	\label{reg:priorita_stesura}
	\begin{enumerate}
		\item \textbf{Vettori obbligati prima di tutto}: sensibile al soprano sempre verso la tonica; settima sempre verso la risoluzione discendente. Questi movimenti non sono negoziabili.
		\item \textbf{Note comuni}: tenere nella stessa voce le note che i due accordi consecutivi hanno in comune.
		\item \textbf{Moto contrario con il basso}: quando possibile, le voci superiori vanno in direzione opposta al basso.
		\item \textbf{Grado congiunto}: le voci restanti si muovono per la distanza minima.
		\item \textbf{Cantabilità del soprano}: scegliere la posizione dell'accordo (ottava, terza, quinta) che produce la linea melodica più fluida e cantabile al soprano.
	\end{enumerate}
\end{regola}

\subsection{La scelta della posizione iniziale}

La prima decisione da prendere è quale nota del primo accordo mettere al soprano. Questo determina la posizione di partenza e influenza tutta la condotta successiva.

\begin{regola}[Scelta della posizione iniziale]
	\label{reg:posizione_iniziale}
	Se non ci sono indicazioni specifiche, la posizione di ottava (fondamentale al soprano) è la scelta più solida e neutra per iniziare. Essa favorisce la cadenza perfetta finale (dove il soprano deve arrivare alla tonica). Se il brano è vivace o aperto, la posizione di quinta può dare un inizio più arioso.
\end{regola}

\subsection{Gestione del soprano nei passaggi scalari}

Nei passaggi in cui il basso sale o scende per gradi (Regola dell'Ottava), il soprano deve anch'esso procedere per gradi congiunti o restare su note del rispettivo accordo. È un errore far saltare il soprano mentre il basso procede per gradi.

\section{Fase 5 — Revisione sistematica}
\label{sec:fase5_basso}

Dopo aver scritto l'intero elaborato, si esegue la revisione secondo la checklist degli errori (\ref{reg:checklist_errori}). Si raccomanda di effettuare la lettura:
\begin{enumerate}
	\item Soprano + Basso (verifica delle quinte/ottave parallele e nascoste tra le voci estreme).
	\item Contralto + Tenore (spesso trascurato, ma fertile di errori nascosti).
	\item Soprano + Contralto, poi Tenore + Basso (le coppie adiacenti).
	\item Tutte le voci insieme (verifica visiva della coerenza generale e dei raddoppi).
\end{enumerate}

\section{Esempio svolto n. 1: basso semplice in Do maggiore}
\label{sec:esempio1}

\subsection{Basso dato}
\musica{
	\new Staff {
		\clef bass \key do \major \time 4/4
		\relative do, {
			do2 re | mi fa | sol1 | do,1
		}
	}
}

\subsection{Analisi preliminare}

\textbf{Armatura:} nessuna alterazione $\rightarrow$ Do maggiore o La minore.
\textbf{Nota iniziale:} Do $\rightarrow$ probabilmente Do maggiore (la nota finale confermerà).
\textbf{Nota finale:} Do (lunga) $\rightarrow$ confermato Do maggiore.
\textbf{Alterazioni cromatiche:} nessuna.
\textbf{Struttura:} 4 misure, metro 4/4 (due accordi per misura suggeriti dal ritmo in minime).

\subsection{Segmentazione}

Cadenza attesa in battuta 3 (metà) e in battuta 4 (finale).
\begin{itemize}
	\item Battuta 3: Sol = V grado $\rightarrow$ possibile cadenza sospesa.
	\item Battuta 4: Do = I grado $\rightarrow$ cadenza perfetta finale.
\end{itemize}

\subsection{Classificazione}

\begin{center}
\begin{tabular}{|l|l|l|l|}
\hline
\textbf{Batt.} & \textbf{Nota basso} & \textbf{Grado} & \textbf{Accordo scelto} \\
\hline
1 & Do (forte) & I & I (stato fond., 5) \\
1 & Re & II, congiunto & VII6 (primo rivolto) \\
2 & Mi & III, congiunto & I6 (primo rivolto) \\
2 & Fa & IV, congiunto & II6 (primo rivolto) \\
3 & Sol (lunga, forte) & V & V (stato fond.) \\
4 & Do (lunga, forte) & I & I (stato fond.) \\
\hline
\end{tabular}
\end{center}

\subsection{Realizzazione}

\musica{
	\partilate
	{ \key do \major \time 4/4 }
	{ do''2 si'2 | do''2 la'2 | si'1 | do''1 }
	{ mi'2 fa'2 | do'2 re'2 | sol'1 | sol'1 }
	{ sol2 fa2 | sol2 la2 | re'1 | mi'1 }
	{ do2 re2 | mi2 fa2 | sol1 | do1 }
	\lyricmode { "I"2 "VII°6" | "I6" "II6" | "V"1 | "I"1 }
	\figuremode { s2 <6> | <6> <6> | s1 | s1 }
}

\subsection{Commento}

La realizzazione applica sistematicamente la Regola dell'Ottava per i gradi congiunti della prima e della seconda battuta. Il soprano forma una linea cantabile Do-Si-Do-La, che sale poi alla sensibile Si (battuta 3) e arriva al Do. La cadenza sospesa in battuta 3 crea la tensione attesa, risolta dalla cadenza perfetta in battuta 4 con il soprano all'ottava (Do'').

\section{Esempio svolto n. 2: basso con progressione e modulazione}
\label{sec:esempio2}

\subsection{Basso dato}
\musica{
	\new Staff {
		\clef bass \key re \minor \time 4/4
		\relative re, {
			re2 mi | fa2 sol | la1 | re,1 \bar "||"
			re2 mi | fa2 fad | sol1 | re1 \bar "|."
		}
	}
}

\subsection{Analisi preliminare}

\textbf{Armatura:} 1 bemolle (Si\bemolle{}) $\rightarrow$ Fa maggiore o Re minore.
\textbf{Nota iniziale e finale:} Re $\rightarrow$ Re minore confermato.
\textbf{Alterazione cromatica:} \Fadiesis{} alla battuta 6 $\rightarrow$ segnale importante da analizzare.
\textbf{Struttura:} 8 misure con doppia stanghetta a metà.

\subsection{Segmentazione}

\begin{itemize}
	\item Prima parte (bb. 1-4): sale da Re a La, cadenza sospesa su La (V grado di Re minore) in battuta 3, poi torna a Re (cadenza perfetta) in battuta 4.
	\item Seconda parte (bb. 5-8): stessa apertura, ma con \Fadiesis{} in battuta 6 $\rightarrow$ tonicizzazione del IV (Sol minore) o modulazione?
	\item La nota finale è Re (4 battute dopo il punto di sospensione) $\rightarrow$ rimaniamo in Re minore. Si tratta quindi di una \textbf{tonicizzazione di Sol minore} (il IV di Re minore), seguita dalla cadenza plagale finale iv--i.
\end{itemize}

\subsection{Classificazione delle note critiche}

\begin{itemize}
	\item \textbf{Fa (battuta 2 e 6)}: III grado di Re minore, armonizzato con i6; il Sol che segue (battuta 2) è il IV grado e prende II°6 (sottodominante). In battuta 6 il \Fadiesis{} (sensibile di Sol) trasforma l'accordo in \textbf{Re7 in primo rivolto}, ovvero V6/5 del IV grado: la progressione Fa-Fa\diesis{}-Sol è un cromatismo ascendente che porta alla tonicizzazione di Sol minore in battuta 7.
\end{itemize}

\subsection{Realizzazione}

\musica{
	\partilate
	{ \key re \minor \time 4/4 }
	{ fa'2 la'2 | la'2 sib'2 | la'1 | la'1 }
	{ re'2 dod'2 | re'2 mi'2 | dod'1 | re'1 }
	{ la2 sol2 | fa2 mi2 | mi1 | fa1 }
	{ re,2 mi, | fa, sol, | la,1 | re,1 }
	\lyricmode { "i"2 "V4/3" | "i6" "II°6" | "V"1 | "i"1 }
	\figuremode { s2 <4 3> | <6> <6> | s1 | s1 }
}

\musica{
	\partilate
	{ \key re \minor \time 4/4 }
	{ fa'2 la'2 | la'2 la'2 | sib'1 | la'1 }
	{ re'2 dod'2 | re'2 re'2 | re'1 | re'1 }
	{ la2 sol2 | fa2 do'2 | sib1 | fa1 }
	{ re,2 mi, | fa, fad, | sol,1 | re,1 }
	\lyricmode { "i"2 "V4/3" | "i6" "V6/5/iv" | "iv"1 | "i"1 }
	\figuremode { s2 <4 3> | <6> <6 5> | s1 | s1 }
}

\subsection{Commento}

La prima parte è una semplice progressione scalare ascendente armonizzata con la Regola dell'Ottava. La seconda parte è identica fino alla battuta 6, dove il \Fadiesis{} trasforma l'accordo sul Fa in V6/5 del IV (Re7 in primo rivolto), creando una spinta cromatica verso Sol minore. La cadenza finale è plagale (iv--i) in Re minore.

\section{Errori tipici nell'armonizzazione del basso senza numeri}
\label{sec:errori_tipici}

\begin{errore}[Armonizzare nota per nota senza visione d'insieme]
	\label{err:nota_per_nota}
	L'errore più comune dei principianti è quello di analizzare il basso nota per nota senza mai alzare lo sguardo sull'insieme. Il risultato è una successione di accordi corretti singolarmente ma senza logica fraseologica: le cadenze appaiono nei posti sbagliati, le progressioni non vengono riconosciute, le modulazioni sembrano accidentali.
	
	\textbf{Rimedio:} fare sempre prima l'analisi strutturale (fase 2) prima di assegnare accordi.
\end{errore}

\begin{errore}[Usare troppi stati fondamentali in successione scalare]
	\label{err:troppi_fondamentali}
	Usare tutti gli accordi allo stato fondamentale in un basso scalare produce un effetto ``pesante'' e crea facilmente errori di parallelismo.
	
	\textbf{Rimedio:} usare sistematicamente la Regola dell'Ottava con i suoi rivolti (6 per i gradi II, III, IV, VI, VII).
\end{errore}

\begin{errore}[Trascurare la sensibile nel modo minore]
	\label{err:sensibile_minore}
	Nel modo minore, è essenziale alzare il VII grado (sensibile) nella dominante e nei rivolti della dominante. Dimenticarsi di farlo produce una dominante minore senza forza propulsiva.
	
	\textbf{Rimedio:} ogni volta che il basso è sul V grado del modo minore, verificare che la sensibile sia alzata nelle voci superiori.
\end{errore}

\begin{errore}[Confondere tonicizzazione e modulazione]
	\label{err:tonicizzazione_modulazione}
	Trattare ogni alterazione cromatica come una modulazione produce brani che cambiano tonalità ogni misura. Al contrario, ignorare alterazioni che segnalano modulazioni reali produce brani che non hanno senso armonico.
	
	\textbf{Rimedio:} cercare sempre la cadenza di conferma. Se c'è una cadenza autentica nella nuova tonalità, è una modulazione; altrimenti, è una tonicizzazione.
\end{errore}

\begin{errore}[Non risolvere le dissonanze]
	\label{err:dissonanze_non_risolte}
	Introdurre settimi accordi o accordi di sesta eccedente senza poi far risolvere correttamente le note dissonanti.
	
	\textbf{Rimedio:} applicare sempre la regola di risoluzione discendente per grado alla settima, e il movimento obbligato della sensibile.
\end{errore}

\section{Strategie per i casi ambigui}
\label{sec:strategie_ambigue}

Non esiste sempre una sola soluzione corretta. Quando il contesto è ambiguo, usare le seguenti \textbf{strategie di scelta}:

\begin{regola}[Strategie per l'ambiguità armonica]
	\label{reg:strategie_ambiguita}
	\begin{itemize}
		\item \textbf{In caso di dubbio, privilegiare la soluzione più semplice}: accordo di tonica o di dominante invece di dominanti secondarie complesse.
		\item \textbf{Preferire la coerenza alla varietà}: se una progressione funziona, continuarla. Non cambiare strategia a metà senza una ragione armonica chiara.
		\item \textbf{Valutare sempre il soprano}: tra due soluzioni armonicamente equivalenti, scegliere quella che produce la linea melodica più cantabile al soprano.
		\item \textbf{Verificare la struttura delle cadenze}: le cadenze devono trovarsi nei punti giusti. Se un'armonizzazione produce due cadenze perfette consecutive, qualcosa non va nella pianificazione fraseologica.
		\item \textbf{La nota del basso ha funzione tonale nel suo contesto}: ricordare che una nota non ha significato da sola. Se un Sol appare dopo un La-Si-Do-Re in scala ascendente di Mi minore, quasi certamente è il III grado (Sol$\natural$) della scala melodica, non il dominante di Do.
	\end{itemize}
\end{regola}

\section{Esercizi graduati}
\label{sec:esercizi_graduati}

I seguenti bassi sono ordinati per difficoltà crescente. Per ciascuno, seguire le cinque fasi del metodo prima di scrivere le voci superiori.

\subsection{Livello 1: basso scalare in modo maggiore (senza alterazioni)}

\musica{
	\new Staff {
		\clef bass \key sol \major \time 4/4
		\relative sol, {
			sol2 la | si do' | re1 | sol,1
		}
	}
}

\textbf{Suggerimenti:} tonalità Sol maggiore; scala ascendente I-II-III-IV, poi V e tonica; applicare Regola dell'Ottava nella prima metà; cadenza sospesa (V) a metà e perfetta alla fine.

\subsection{Livello 2: basso misto (congiunto e salti) in modo minore}

\musica{
	\new Staff {
		\clef bass \key la \minor \time 4/4
		\relative la, {
			la2 la | re2 mi | fa2 mi | la,1
		}
	}
}

\textbf{Suggerimenti:} La minore; battuta 1, La ripetuta $\rightarrow$ I, possibile cambio posizione; battuta 2, Re-Mi $\rightarrow$ IV-V; battuta 3, Fa-Mi $\rightarrow$ VI-V; battuta 4, cadenza perfetta. Attenzione alla sensibile (\Soldiesis{}) nel V grado.

\subsection{Livello 3: basso con progressione per quinte discendenti}

\musica{
	\new Staff {
		\clef bass \key do \major \time 3/4
		\relative do, {
			do4 fa si, | mi la re | sol do2 | do1.
		}
	}
}

\textbf{Suggerimenti:} Do maggiore; le prime due battute formano una progressione per quinte discendenti; alternare stati fondamentali e rivolti; battuta 3 è cadenza composta IV-V-I (Do-Sol); battuta 4 è la tonica finale.

\subsection{Livello 4: basso con tonicizzazione}

\musica{
	\new Staff {
		\clef bass \key mi \minor \time 4/4
		\relative mi, {
			mi2 fad | sol la | si2 lad | si1
		}
	}
}

\textbf{Suggerimenti:} Mi minore (1 diesis); in battuta 3, il \Ladiesis{} è la sensibile di Si (V di Mi m) $\rightarrow$ tonicizzazione del V; il brano inizia e finisce in Mi minore; cadenza sospesa su Si in battuta 4 è in realtà una cadenza perfetta nella tonicizzazione di Si, ma vista in Mi minore è una cadenza sospesa.

\subsection{Livello 5: basso con modulazione a un tono vicino}

\musica{
	\new Staff {
		\clef bass \key fa \major \time 4/4
		\relative fa, {
			fa2 sol | la sib | do2 re | sib1 \bar "||"
			do2 re | mi fa | sol1 | do1 \bar "|."
		}
	}
}

\textbf{Suggerimenti:} Fa maggiore (prima parte); il Mi$\natural$ nella seconda parte (battuta 5) è la sensibile di Fa? No --- in Fa il Mi è già la sensibile naturale. Allora, il Do della seconda parte potrebbe essere il I di Do maggiore. Verificare: la seconda parte finisce su Do (tonica di Do maggiore) $\rightarrow$ modulazione a Do maggiore (tonalità della dominante di Fa). L'accordo pivot è il Sol (II di Fa M = III di Sol M oppure V di Do M) --- individuarlo e usarlo come cerniera.

\section{Dalla procedura all'intuizione}
\label{sec:dall_procedura_intuizione}

Il metodo descritto in questo capitolo è, per sua natura, una procedura esplicita e sistematica. È così che deve essere per chi si avvicina per la prima volta al basso senza numeri: un insieme di passi chiari, controllabili, verificabili.

L'obiettivo a lungo termine, tuttavia, è la \textbf{fluidità}: la capacità di leggere un basso e sentire immediatamente le funzioni armoniche, senza doverle dedurre passo per passo. Questa fluidità si acquisisce con la pratica ripetuta e riflessiva --- non con la memorizzazione delle regole, ma con la loro \textit{applicazione consapevole} su molti bassi diversi.

\begin{osservazione}[Il ruolo dell'analisi dei brani]
	\label{oss:analisi_brani}
	Un acceleratore potente dello sviluppo dell'intuizione armonica è l'\textbf{analisi dei brani reali}: corali di Bach, sonate classiche, mottetti rinascimentali. Leggere un basso analizzato (dove si sa già qual è l'armonia) e chiedersi ``come avrei dedotto questo accordo?'' è un esercizio che sviluppa la capacità di riconoscimento in modo molto più efficace di qualsiasi schema astratto.
\end{osservazione}

L'armonizzazione del basso senza numeri è la dimostrazione pratica che l'armonia è stata compresa, non solo memorizzata. Non è la fine del percorso, ma l'inizio della vera autonomia musicale.
